\documentclass[11pt]{article}

\usepackage[margin=1in]{geometry}
\usepackage{amsmath,amssymb,amsthm}
\usepackage{booktabs}
\usepackage{array}
\usepackage{longtable}
\usepackage{xcolor}
\usepackage[hidelinks]{hyperref}

\DeclareMathOperator*{\argmin}{arg\,min}
\DeclareMathOperator*{\argmax}{arg\,max}
\newcommand{\R}{\mathbb{R}}
\newcommand{\Dstar}{\mathcal{D}^*}

\title{DUST and DUSTib Pruning in the \texttt{dust} Package}
\author{Code and mathematical implementation report}
\date{October 2026}

\begin{document}

\maketitle

\begin{abstract}
This report compares the two one-constraint pruning policies called
\texttt{DUST} and \texttt{DUSTib} in the current \texttt{dust} package.
The comparison is made against the derivation in the DUST paper, especially
the one-dimensional decision function and its closed-form maximizer.  The
main conclusion is that the two policies are intended to implement the same
mathematical pruning test.  They are not different statistical methods.
\texttt{DUST} is a more optimized implementation of the same maximum, whereas
\texttt{DUSTib} evaluates the decision-function maximizer more literally.
The current \texttt{DUSTib} code has boundary and stationary-point branches
that can cause under-pruning, and it is not used by the Highway SIMD path.
The report gives a minimal correction, a more robust design, and a testing
plan for making both approaches exact, effective, and fast.
\end{abstract}

\tableofcontents

\section{Scope and sources}

The mathematical reference used here is the local DUST manuscript
\texttt{PAPER\_DUST/DUST.tex}.  The relevant material is the section
``DUST method with one constraint'', the proof of the closed-form maximum in
the appendix, and the appendix tables describing the convex conjugates and
their domains.  In particular, the paper defines the decision function in
Equation~\texttt{eq:decisionDual1D} and derives its maximizer in the proof of
Theorem~\texttt{prop:dual1D\_max\_exact}.

The implementation examined in this report is:

\begin{itemize}
  \item \texttt{dust/src/1D\_DualMaxPolicies.h}, which contains
        \texttt{DualMax\_DUST} and \texttt{DualMax\_DUSTib};
  \item \texttt{dust/src/1D\_DUST\_Impl.h}, which computes the dynamic
        program and calls the selected pruning policy;
  \item \texttt{dust/src/DUST\_1D\_HW.cpp}, which provides the Highway SIMD
        implementation and explicitly delegates \texttt{DUSTib} to the
        scalar engine;
  \item the model policies
         \texttt{1D\_[1--8]\_*.h}, which provide $\Dstar$, its derivatives,
        domains, boundary limits, and segment costs.
\end{itemize}

The paper does not use the name \texttt{DUSTib}.  That name belongs to an
implementation variant inherited from the older \texttt{dust} code.  It is
therefore important not to interpret \texttt{DUSTib} as a second theoretical
pruning method introduced by the paper.

\section{Common dynamic-programming problem}

Let $y_{st}$ denote the observations $y_{s+1},\ldots,y_t$, with
$0\leq s<t\leq n$.  For a one-parameter natural exponential family, the
segment cost has the canonical form
\begin{equation}
  c(y_{st};\theta)=(t-s)A(\theta)-\theta S_{st},
  \qquad S_{st}=\sum_{i=s+1}^t T(y_i),
  \label{eq:canonical-cost}
\end{equation}
where $A$ is the log-partition function and $T$ is the sufficient statistic.
The dynamic-programming cost is
\begin{equation}
  Q_t=\min_{0\leq s<t}\{Q_s+c(y_{st})+\beta\},
  \qquad Q_0=-\beta,
  \label{eq:op}
\end{equation}
where $c(y_{st})=\min_\theta c(y_{st};\theta)$ and $\beta$ is the segment
penalty in the convention used by the package.

For each possible last change $s$, define the function-valued cost
\begin{equation}
  q_t^s(\theta)=Q_s+c(y_{st};\theta)+\beta.
  \label{eq:functional-cost}
\end{equation}
The purpose of pruning is to remove an index $s$ when it can never be the
last change of an optimal future segmentation.

\subsection{The one-constraint pruning problem}

Suppose that $r<s$.  The one-constraint relaxation used by DUST compares
$q_t^s$ with $q_t^r$:
\begin{equation}
\left\{
\begin{aligned}
  &\min_{\theta\in\Theta}q_t^s(\theta),\\
  &\text{s.t.}\quad q_t^s(\theta)-q_t^r(\theta)\leq 0.
\end{aligned}
\right.
\label{eq:pruning-problem}
\end{equation}
The convex conjugate associated with $A$ is
\begin{equation}
  \Dstar(x)=x(\nabla A)^{-1}(x)-A\bigl((\nabla A)^{-1}(x)\bigr).
  \label{eq:conjugate}
\end{equation}
The paper uses the following four scalar quantities, which correspond
directly to the variables named $a,b,c,d$ in the C++ code:
\begin{align}
  a&=\overline S_{st}=\frac{S_{st}}{t-s},
  &b&=\overline S_{rs}=\frac{S_{rs}}{s-r},
  \label{eq:ab}\\
  c&=\overline Q_{st}=\frac{Q_t-Q_s}{t-s},
  &d&=\overline Q_{rs}=\frac{Q_s-Q_r}{s-r}.
  \label{eq:cd}
\end{align}

\section{The decision function and its maximum}

After changing variables from the nonnegative Lagrange multiplier to
$x=\mu/(1-\mu)$, the paper obtains the decision function
\begin{equation}
  \mathbb{D}(x)
  =-\Dstar\bigl(a+x(a-b)\bigr)-\bigl(c+x(c-d)\bigr),
  \qquad x\in[0,x_{\max}).
  \label{eq:decision}
\end{equation}
The test is safe because a positive value proves that the constrained
minimum is above the pruning threshold.  Thus the pruning condition is
\begin{equation}
  \max_{0\leq x<x_{\max}}\mathbb{D}(x)>0.
  \label{eq:prune}
\end{equation}

The function is concave because $\Dstar$ is convex.  If $a\neq b$, its
stationary point satisfies
\begin{align}
  0&=-(a-b)(\Dstar)'\bigl(a+x^\star(a-b)\bigr)-(c-d),
  \\[-2pt]
  R&=-\frac{c-d}{a-b},
  \label{eq:R}\\
  m^\star&=((\Dstar)')^{-1}(R),
  \\[-2pt]
  x^\star&=\frac{m^\star-a}{a-b}.
  \label{eq:xstar}
\end{align}
In the C++ policies, \texttt{R} is the natural-parameter value in the
stationarity equation.  The model call \texttt{DstarPrimeInv}$(R)$ returns
$m^\star$.
It is a mistake to compare $R$ directly with a boundary expressed in the
mean-parameter scale.  The current \texttt{DUST} code explicitly avoids that
mistake by comparing $\texttt{DstarPrimeInv(R)}$ with the mean at the boundary.

The maximizer is selected by the following rules:

\begin{enumerate}
  \item If $x^\star\leq0$, the maximum on the feasible domain is at $x=0$.
        This is exactly the PELT test:
        \begin{equation}
          \mathbb{D}(0)=-\Dstar(a)-c.
          \label{eq:pelt}
        \end{equation}
  \item If $0<x^\star<x_{\max}$, the maximum is the interior value
        $\mathbb{D}(x^\star)$.
  \item If $x^\star\geq x_{\max}$ and $x_{\max}$ is finite, the supremum is
        the limiting value at the endpoint.  If $a>b$, the path moves toward
        the right boundary of the mean domain; if $a<b$, it moves toward the
        left boundary.
  \item If $x_{\max}=+\infty$, the correct treatment depends on the
        asymptotic derivative and, in equality cases, on the asymptotic value
        of the decision function.  This case must not be confused with a
        finite endpoint.
\end{enumerate}

The paper states the essential negative-stationary-point rule explicitly:
when $x^\star<0$, the decision function is evaluated at zero and the test is
the PELT inequality.  There is no special mathematical threshold at
$x^\star=-1$.

\subsection{Model domains and endpoint directions}

The model-specific conjugates used in the paper and the code are summarized
below.  The table is also the reason endpoint handling cannot be hard-coded
to the left boundary.

\begin{center}
\small
\begin{tabular}{lll}
\toprule
Model & $\Dstar(m)$ & Mean-domain information \\
\midrule
Gaussian & $m^2/2$ & $m\in\R$ \\
Exponential & $-\log(m)-1$ & $m>0$ \\
Poisson & $m(\log m-1)$ & $m>0$ \\
Geometric & $(m-1)\log(m-1)-m\log m$ & $m>1$ \\
Bernoulli/Binomial & $m\log m+(1-m)\log(1-m)$ & $0<m<1$ \\
Negative Binomial & $m\log m-(1+m)\log(1+m)$ & $m>0$ \\
Variance & $-\tfrac12(\log m+1)$ & $m>0$ \\
\bottomrule
\end{tabular}
\end{center}

For Bernoulli and Binomial, the domain is bounded on both sides.  Therefore,
when $a>b$, $x_{\max}$ is a right-boundary value and the endpoint must use
$\Dstar(1)$, not $\Dstar(0)$.  For all other models, finite $x_{\max}$ in
the $a<b$ direction reaches the left boundary.  The paper's Table
\texttt{dual\_table} gives the corresponding formulas for $x_{\max}$.

\section{What DUST and DUSTib mean in the code}

\subsection{\texttt{DualMax\_DUST}}

The current \texttt{DualMax\_DUST} implementation in
\texttt{1D\_DualMaxPolicies.h:76--157} is an optimized closed-form test.
Its main steps are:

\begin{enumerate}
  \item compute $a,b,c,d$ and the model-dependent domain limit;
  \item handle a boundary value or a linear decision function;
  \item evaluate the derivative at $x=0$; if the derivative is negative,
        return the PELT test immediately;
  \item handle an infinite domain using the derivative at infinity;
  \item for a finite domain, compute the actual critical mean
        $(\Dstar)'^{-1}(R)$ and compare it with the mean at the endpoint;
  \item for an interior maximum, use the segment-cost shortcut
        \texttt{Model::costEval(..., R, ...)} instead of explicitly
        evaluating the full decision function at $x^\star$.
\end{enumerate}

The final shortcut is mathematically equivalent to evaluating
$\mathbb{D}(x^\star)$.  It follows from the Legendre relation and the
stationarity equation, and it avoids constructing $x^\star$ and the
interior mean when the critical point is known to be feasible.

\subsection{\texttt{DualMax\_DUSTib}}

The current \texttt{DualMax\_DUSTib} implementation in
\texttt{1D\_DualMaxPolicies.h:159--202} takes a more literal route:

\begin{enumerate}
  \item compute $R$;
  \item compute $m^\star=((\Dstar)')^{-1}(R)$ and
        $x^\star=(m^\star-a)/(a-b)$;
  \item if the point is considered interior, evaluate
        \begin{equation}
          -\Dstar\bigl(a+x^\star(a-b)\bigr)
          -\bigl(c+x^\star(c-d)\bigr);
          \label{eq:dustib-direct}
        \end{equation}
  \item otherwise use a PELT or boundary/asymptotic branch.
\end{enumerate}

The intended formula is therefore the same as \texttt{DUST}.  The difference
is implementation strategy, not a different pruning criterion:

\begin{center}
\small
\begin{tabular}{p{0.25\linewidth}p{0.31\linewidth}p{0.31\linewidth}}
\toprule
Property & \texttt{DUST} & \texttt{DUSTib} \\
\midrule
Mathematical target & Maximum of Equation~\eqref{eq:decision} & Same maximum \\
Interior evaluation & Legendre/segment-cost shortcut & Direct decision-function value \\
Early PELT branch & Checks the derivative at $x=0$ first & Computes the inverse derivative before its PELT branch \\
Endpoint handling & Selects left or right boundary from the direction of $a-b$ & Current code uses the left boundary in both finite-endpoint branches \\
Highway path & Hand-vectorized for DUST & Explicitly falls back to the scalar engine \\
Status in documentation & Default and recommended exact rule & Marked ``not yet working'' in the R documentation \\
\bottomrule
\end{tabular}
\end{center}

\section{Why DUST is faster}

There are four separate reasons.  They should not be conflated.

\subsection{The main reason: SIMD is unavailable for DUSTib}

The file \texttt{DUST\_1D\_HW.cpp:53--62} says that \texttt{DUSTib} is not
hand-vectorized and delegates to the scalar \texttt{dust} engine.  The
dispatch at lines 109--115 confirms that \texttt{DUSTib} calls
\texttt{run\_dust\_engine<Model, DualMax\_DUSTib<Model>>}.  In contrast,
\texttt{DUST} uses the Highway implementation when the package is built with
Highway.  Its active candidates are stored in a flat array so that candidate
values can be loaded in SIMD lanes.  The scalar engine uses the linked-list
style active-index structure in \texttt{1D\_Indices.h}; this structure is not
directly vectorizable.

Consequently, comparing \texttt{DUST.1D.HW} with method \texttt{DUST} to
the same entry point with method \texttt{DUSTib} is primarily a
SIMD-versus-scalar comparison as well as a pruning-rule comparison.

\subsection{DUST does less work per surviving candidate}

The DUST policy checks the derivative at zero before calculating the inverse
derivative.  When the derivative is negative, the maximum is already known to
be the PELT point, so the expensive inverse and endpoint calculations are
skipped.  The DUSTib policy currently computes
\texttt{DstarPrimeInv(R)} and $x^\star$ before making its negative-point
decision.

For an interior maximum, DUST uses \texttt{costEval}, which is a compact
model-specific expression for the value at the stationary point.  DUSTib
constructs the stationary mean and evaluates $\Dstar$ plus the affine term.
Both are constant-time operations, but the latter generally performs more
floating-point work and more domain-sensitive logarithm evaluations.

\subsection{Under-pruning increases later dynamic-programming work}

Any missed positive decision-function value leaves a candidate in the active
set.  The next iteration then evaluates that candidate in the OP step, and it
can be tested again at later times.  Thus a branch error in DUSTib does not
change the exact segmentation, provided it only causes under-pruning, but it
can substantially increase runtime.  The current DUSTib negative-point and
endpoint branches can cause precisely this type of under-pruning.

\subsection{The comparison must be benchmarked fairly}

On a scalar build, or when comparing \texttt{dust.1D} rather than
\texttt{DUST.1D.HW}, the SIMD explanation does not apply.  The remaining
difference is then due to operation count, branch behavior, active-set size,
and numerical branch outcomes.  Benchmarks should report separately:

\begin{itemize}
  \item scalar DUST versus scalar DUSTib;
  \item Highway DUST versus scalar DUSTib;
  \item final active-set size and the complete active-set trajectory;
  \item identical data, penalty, model normalization, compiler, and hardware.
\end{itemize}

\section{DUSTib corrections needed for full effectiveness}

\subsection{\texorpdfstring{Correction 1: every negative $x^\star$ is a PELT case}{Correction 1: every negative xstar is a PELT case}}

The current code contains:

\begin{verbatim}
if(xstar <= 0 && xstar > -1)
{
  return (- Model::Dstar(a) -c > 0);
}
\end{verbatim}

The lower bound $-1$ has no role in the decision-function domain.  The
change of variables used in the paper gives $x\geq0$, so any $x^\star<0$ is
outside the feasible domain, regardless of how negative it is.  The literal
correction is:

\begin{verbatim}
if (xstar <= 0.0)
{
  return (-Model::Dstar(a) - c > 0);
}
\end{verbatim}

This is exactly the statement in the paper's proof: when $x^\star<0$, the
maximum is at zero and the test is PELT.  The same rule should be used with a
small numerical tolerance, for example $x^\star\leq\varepsilon_x$.

\subsection{Correction 2: use the boundary reached by the path}

The current DUSTib code uses \texttt{Dstar\_leftboundary()} in both finite
endpoint branches:

\begin{verbatim}
return (-Model::Dstar_leftboundary()
        - (c + x_max * (c-d))) > 0;
\end{verbatim}

This is wrong for a bounded-right-domain model when $a>b$.  Since
$m(x)=a+x(a-b)$, the path moves upward when $a>b$ and reaches the right
boundary.  The endpoint value must therefore be selected as:

\begin{verbatim}
  ? Model::Dstar_rightboundary()
  : Model::Dstar_leftboundary();

return (-Dstar_at_xmax
        - (c + x_max * (c-d))) > 0;
\end{verbatim}

This matters for Bernoulli and Binomial, whose mean domain is $(0,1)$ and
whose right boundary is $1$.  The current \texttt{DUST} policy already uses
this directional selection at lines 146--152; DUSTib should use the same
logic.

\subsection{Correction 3: do not treat an infinite domain as a finite endpoint}

The current test

\begin{verbatim}
if(xstar >= x_max) // case finite only
\end{verbatim}

is not actually restricted to a finite $x_{\max}$.  If both values are
infinite, the expression can enter a branch containing
$x_{\max}(c-d)$, which may be infinite or undefined.  The finite endpoint
branch should explicitly require:

\begin{verbatim}
if (std::isfinite(x_max) && xstar >= x_max)
{
  // evaluate the directional boundary limit
}
\end{verbatim}

The $x_{\max}=+\infty$ case should instead be handled by a model policy that
provides the derivative limit at infinity and, when necessary, the value of
the decision function at infinity.

\subsection{Correction 4: make out-of-domain stationary points explicit}

For Exponential, Geometric, Negative Binomial, and Variance models, the
natural parameter has a one-sided domain.  The inverse derivative
\texttt{DstarPrimeInv(R)} is meaningful only when $R$ belongs to the range of
\texttt{DstarPrime}.  If it returns NaN, an invalid mean, or an infinite mean,
the code must not silently continue to the generic asymptotic branch.

A robust policy should expose a model-level operation that returns either a
valid critical mean or an ``out of range'' result.  An out-of-range critical
point is then handled by the appropriate endpoint rule.  A safe numerical
fallback is to return ``do not prune'' if the endpoint cannot be evaluated
reliably; this sacrifices speed but preserves exactness.

\subsection{Correction 5: replace the ambiguous asymptotic helper}

The scalar policies expose \texttt{Dstar\_superLinearLimit()}, but the
DUSTib final branch uses it as if it were enough to decide the behavior at an
unbounded endpoint.  Its name and values do not consistently represent the
limit of $\Dstar$.  For example, for the Exponential model,
$\Dstar(m)=-\log(m)-1\to-\infty$ as $m\to\infty$, while
$(\Dstar)'(m)\to0$.

The model interface should distinguish at least:

\begin{verbatim}
DstarPrimeAtInfinity()       // limit of the derivative
dualAtInfinity(c, d, a, b)   // sign/value of the full endpoint limit
\end{verbatim}

The first quantity is enough when the asymptotic slope is strictly positive
or negative.  The second is needed in equality cases, such as an asymptotic
slope of zero.  The current DUST policy's use of
\texttt{DstarPrime(+infinity)} is closer to the correct first-order test, but
it should also document and handle equality cases explicitly.

\subsection{A corrected branch structure}

The following is a schematic replacement for the DUSTib branch logic.  It is
not a complete patch because the scalar model interface first needs a
well-defined invalid-domain and asymptotic result, but it shows the correct
ordering:

\begin{verbatim}
// a, b, c, d already computed; require a != b here.
const double f0 = -Model::Dstar(a) - c;
const double g0 = -(a-b) * Model::DstarPrime(a) - (c-d);

// The feasible decision-function domain starts at x = 0.
if (g0 <= 0.0) return f0 > 0.0;

const double R = -(c-d) / (a-b);
const CriticalMean critical = Model::criticalMean(R);

if (!critical.valid) {
  // The stationary point is outside the mean domain.
  // Select the endpoint from the sign of the derivative/domain direction.
  return endpointOrAsymptoticTest<Model>(a, b, c, d);
}

const double xstar = (critical.mean - a) / (a-b);
const double xmax = Model::xMax(a, b);

if (xstar <= 0.0) return f0 > 0.0;

if (std::isfinite(xmax) && xstar >= xmax)
  return directionalEndpointTest<Model>(a, b, c, d, xmax);

if (std::isfinite(xstar) && xstar < xmax)
  return decisionAt<Model>(a, b, c, d, xstar) > 0.0;

return asymptoticTest<Model>(a, b, c, d);
\end{verbatim}

At minimum, the first two literal corrections should be applied immediately:
replace the $-1$ condition by $x^\star\leq0$ and select the endpoint boundary
according to the sign of $a-b$.  The remaining changes make the method robust
at domain and infinity boundaries rather than merely fixing the common
interior cases.

\section{Improvements for both approaches}

\subsection{Use one mathematically authoritative maximization kernel}

The two policies currently duplicate delicate model-dependent branch logic.
That makes it easy for one policy to receive a boundary fix while the other
does not.  A better design is to centralize the following operations:

\begin{itemize}
  \item the domain of the mean statistic;
  \item the map $(\Dstar)'$ and its inverse;
  \item left and right boundary values;
  \item the directional endpoint selected by $a-b$;
  \item finite and infinite endpoint evaluation;
  \item a checked stationary-point result.
\end{itemize}

Then DUST and DUSTib can differ only in how they evaluate a valid interior
maximum.  DUST can retain the segment-cost shortcut; DUSTib can retain the
direct decision-function formula for diagnostic purposes.  Both must call
the same endpoint and invalid-domain helpers.

\subsection{Make DUSTib a debug/reference mode, not a second unchecked kernel}

Once the shared maximization logic is available, DUSTib is useful as a
reference implementation.  For every candidate, it can evaluate the direct
decision function and compare its result with the DUST shortcut in debug
builds.  This provides a strong regression check for the Legendre shortcut
and for model-specific formulas without imposing the extra work on the
production path.

\subsection{Vectorize only after correctness is established}

If DUSTib is retained as a user-facing method, it should either use the same
Highway lane kernel as DUST or be documented as scalar reference code.  A
second hand-vectorized branch tree would increase maintenance cost.  The
preferred performance path is to vectorize the shared decision test and use
lane masks for the PELT, interior, endpoint, and asymptotic cases.

\subsection{Strengthen numerical safeguards}

The implementation should use model-specific stable primitives and checked
domains throughout:

\begin{itemize}
  \item use \texttt{log1p} and stable logistic evaluations near $0$ and $1$;
  \item avoid evaluating $\log(0)$ or $0\log(0)$, even when a floating-point
        average is within a few ulps of a boundary;
  \item use a named tolerance policy rather than scattered constants such as
        \texttt{1e-9} and \texttt{1e-14};
  \item distinguish a finite endpoint, an infinite endpoint, and an invalid
        stationary point;
  \item never prune when the decision value is NaN or the domain status is
        uncertain.
\end{itemize}

The last rule is conservative: it may reduce pruning, but it cannot destroy
the exactness guarantee.

\subsection{Test pruning decisions, not only final segmentations}

The current tests in \texttt{tests/testthat/test-dust1D.R} verify that the
methods return the same final change points on selected examples.  That is
necessary but insufficient.  Different pruning schedules can produce the
same final segmentation even when one method misses many valid pruning
opportunities.

The test suite should additionally compare each policy with an independent
reference maximization of Equation~\eqref{eq:decision} on small examples.
The reference should test:

\begin{enumerate}
  \item $x^\star$ slightly below zero and far below zero;
  \item an interior $x^\star$;
  \item $x^\star$ beyond a finite left endpoint;
  \item $x^\star$ beyond a finite right endpoint for Bernoulli/Binomial;
  \item $x_{\max}=+\infty$;
  \item $a=b$, $a$ at a model boundary, and $a$ numerically close to a
        boundary;
  \item invalid or non-finite inverse-derivative values.
\end{enumerate}

The reference should compare the Boolean pruning decision and the active-set
trajectory, not merely $Q_t$ and the backtracked change points.  For small
$n$, an exhaustive grid plus a high-accuracy one-dimensional optimizer is
adequate.  The production policy should never prune when the reference says
the maximum is non-positive.

\subsection{Improve candidate selection for higher-dimensional extensions}

The paper recommends the nearest smaller active index as an effective
single-constraint choice, which is what the current one-dimensional code
uses.  This is a good default for one-parameter models.  For multi-parameter
models, stronger lower bounds can be obtained by using several constraints,
up to the parameter dimension under the paper's strong-duality result.
The implementation should therefore separate:

\begin{itemize}
  \item the candidate-selection policy;
  \item the number of constraints;
  \item the maximizer for the resulting one- or multi-dimensional decision
        function.
\end{itemize}

This avoids treating a one-dimensional closed form as if it automatically
solved the multi-parameter problem.

\subsection{Keep benchmarks comparable}

The paper's conclusions concern exact pruning and measured computation time,
not just arithmetic complexity.  Future comparisons should record the model
normalization, penalty convention, active-set size, hardware, compiler,
scalar/SIMD backend, and whether the reference method is allowed to apply a
smallest-index PELT check.  In particular, the current
\texttt{dust.object.1D.HW} documentation correctly notes that
\texttt{DUSTib} falls back to the scalar object; this fact must be visible in
any speed table.

\section{Conclusions}

The DUST paper gives one one-constraint decision function and one closed-form
maximum for the one-parameter case.  The repository names two implementations
of that idea \texttt{DUST} and \texttt{DUSTib}; they should have the same
mathematical result.  \texttt{DUST} is faster because it performs an early
PELT check, uses a Legendre/segment-cost shortcut, and, in the Highway entry
point, is vectorized.  \texttt{DUSTib} currently performs more scalar work and
falls back to the scalar engine.

The most important DUSTib correction is to replace the condition
$-1<x^\star\leq0$ by the mathematically correct condition $x^\star\leq0$.
The next essential correction is to use the right boundary when $a>b$,
particularly for Bernoulli and Binomial models.  Infinite endpoints and
out-of-domain stationary points should then be handled through an explicit
model-policy interface rather than through the current ambiguous final
branch.

After these corrections, the two policies should share a tested maximization
kernel.  DUST can remain the optimized production method, while DUSTib can be
retained as a transparent reference/debug implementation or vectorized if a
separate production option is genuinely needed.

\begin{thebibliography}{9}

\bibitem{dustpaper}
V. Runge, C. Truong, and S. Quern\'{e}.
\newblock DUST: A Duality-Based Pruning Method For Exact Multiple Change-Point Detection.
\newblock Local source: \texttt{PAPER\_DUST/DUST.tex}, especially the
decision-function equation and the proof of the one-dimensional maximum.

\bibitem{dustcode}
\texttt{dust} source code.
\newblock In particular: \texttt{src/1D\_DualMaxPolicies.h},
\texttt{src/1D\_DUST\_Impl.h}, and \texttt{src/DUST\_1D\_HW.cpp}.

\end{thebibliography}

\end{document}
